\begin{textsample}{POS Dim 3 – es – Score 19.00 – es/es\_inf\_7.txt}  \label{ex:f3_pos_118}
Os Campos de Experiências e os Objetivos de Aprendizagem e Desenvolvimento Assim, a organização curricular da Educação \textbf{Infantil} estrutura -se em cinco Campos de Experiências, no âmbito dos quais são definidos os Objetivos de Aprendizagem e Desenvolvimento : - O Eu, O Outro e o Nós - Corpo, \textbf{Gestos} e Movimentos - Traços, \textbf{Sons}, Cores e Formas - \textbf{Escuta}, Fala, Pensamento e Imaginação - Espaços, Tempos, Quantidades, Relações e Transformações Figura 18 – Os Campos de Experiência da Educação \textbf{Infantil} Os Objetivos de Aprendizagem e Desenvolvimento, que são específicos para o currículo da Educação \textbf{Infantil}, ao serem pensados por cada \textbf{instituição} escolar, devem detalhar noções, habilidades, atitudes e/ou especificidades locais para cada um dos objetivos de aprendizagem e desenvolvimento da BNCC, além de serem pensados com base nas propostas pedagógicas de cada \textbf{instituição}, que se concretizarão por meio das práticas educacionais organizadas em torno do conhecimento e em meio às relações sociais que se entravam diariamente nesses espaços e afetam a construção das identidades da \textbf{criança}.
Com base nesta concepção, todo planejamento curricular deve ser organizado pensando as diferentes situações que perpassam nos espaços da escola, garantindo à \textbf{criança} seu direito por uma aprendizagem qualitativa, tomando -a como centro em suas decisões, considerando suas linguagens, modo de expressar -se, conhecer, \textbf{desejos}, \textbf{afetos}, abolindo “ todos os procedimentos que não reconheçam a atividade criadora e o protagonismo da \textbf{criança} \textbf{pequena} ou que promovam atividades mecânicas e não significativas para elas ”. ( Parecer CNE/CEB nº 20/09 ).
Embora sejam diversos os posicionamentos relacionados à organização dos grupos \textbf{etários} constitutivos nesta etapa, com uma variação dos \textbf{meses} de um período para o outro, a BNCC situou em três faixas \textbf{etárias} os períodos de desenvolvimento dos \textbf{bebês}, das \textbf{crianças} muito \textbf{pequenas} e das \textbf{crianças} \textbf{pequenas}.
Todavia, esses agrupamentos e suas faixas \textbf{etárias} não podem ser considerados de forma rígida, em virtude dos diferentes ritmos de desenvolvimento, das peculiaridades de cada \textbf{criança} e da forma como cada município se organiza no atendimento a esta etapa de ensino.
Estes grupos \textbf{etários} estão divididos em \textbf{Creche} e Pré-escola, sendo a \textbf{Creche} subdividida em \textbf{Bebês} 0 a 1 ano e 6 \textbf{meses}, \textbf{Crianças} bem \textbf{pequenas} 1 ano e 7 \textbf{meses} a 3 anos e 11 \textbf{meses} e a Pré-escola em \textbf{Crianças} Pequenas – 4 anos a 5 anos e 11 \textbf{meses}, ( BRASIL, 2017, pg. 39 ).
De um modo geral, a LDB Nº 9394/96 estabelece que as etapas \textbf{etárias} da Educação \textbf{Infantil} sejam organizadas : I – \textbf{creches}, ou entidades equivalentes para \textbf{crianças} de zero a três anos de idade ; e II – pré-escolas, para \textbf{crianças} de quatro e cinco anos de idade. ( BRASIL, 1996, art. 30º ) Os Campos de Experiências propõem, no processo de construção de habilidades/competências, uma progressão horizontal e vertical dos objetivos, que assegurem um conjunto de aprendizagens e de desenvolvimento, como direitos essenciais de \textbf{estímulo} à autopercepção, às interações e \textbf{brincadeiras} que proporcionem experiências de aprendizagens significativas e variadas aos \textbf{bebês} e às \textbf{crianças} \textbf{pequenas} e bem \textbf{pequenas}.
O organograma abaixo caracteriza a estrutura da proposta contida no currículo do Espírito Santo, a qual destaca aspectos básicos do desenvolvimento \textbf{infantil}, que devem ser observados e ampliados a partir das práticas pedagógicas locais.
Todos os Campos de Experiências fundamentam -se nos Princípios e nos Direitos da Aprendizagem, tendo como eixos norteadores as interações e \textbf{brincadeiras}, sendo estes elementos básicos na construção de cada \textbf{criança} como ser único, além de serem formas privilegiadas para ela ampliar seus \textbf{afetos}, sensações, \textbf{percepções}, memória, linguagem e formação de sua identidade.
Todo currículo deve se efetivar com base nesses dois processos.
Quanto aos objetivos de aprendizagem e desenvolvimento, alguns encontram maior identidade e possibilidade de serem intencionalmente trabalhados num campo de experiências do que em outros.
Alguns valores e princípios que estes campos encerram são comuns aos demais campos e poderão ser trabalhados, através de diferentes estratégias de ensino, que facilitam a sua incorporação pelas \textbf{crianças}.
A intenção é pensar em ações educativas que contemplem os objetivos da aprendizagem levando as \textbf{crianças} a vivenciarem as experiências adequadas ao seu entendimento.
As estratégias que o professor vai adotar para o efetivo trabalho com os objetivos propostos para as \textbf{crianças} poderão ser variadas, uma vez que elas se estabelecem pelas \textbf{brincadeiras} e interações com seus \textbf{pares} e \textbf{adultos} no cotidiano escolar.
Todos os Campos de Experiências trazem objetivos de aprendizagem e desenvolvimento que devem ser observados em cada um dos grupos \textbf{etários} a serem trabalhados.
Assim, cada um dos cinco campos de experiência é apresentado por tabelas, conforme exemplo abaixo : Entendemos que a tabela não deve ser vista como um roteiro a ser seguido, mas como um conjunto de dicas para orientação do trabalho pedagógico, pois o/a professor/a, por meio da pesquisa, trocas de experiências, planejamento, observação e avaliação das \textbf{crianças} tem condições de evidenciar muitas outras formas trabalho com a \textbf{criança}.

% matched lemmas: adulto, afeto, bebê, brincadeira, creche, criança, desejo, escuta, estímulo, etário, gesto, infantil, instituição, mês, par, pequeno, percepção, som
\end{textsample}
